\chapter{绪论}

\section{思考题}

% source: docs/review_questions.md | chapter=1. 绪论 | question=1
\begin{problem}\label{exer:think_01_intro_01}
请简述新能源与可再生能源的主要区别？
\end{problem}

\begin{solution}
\begin{itemize}
\item \textbf{新能源}：是指在新技术基础上能够开发和利用的能源 。
\item \textbf{可再生能源}：是指在生态循环中能不断再生的能源,不会随着本身的转化或人类的利用而减少,具有天然的自我恢复功能 。
\end{itemize}
常规能源与新能源的划分并无严格界限,会随着时间推移或地区不同发生一定的变化；不可再生能源是指会随着人类的使用而越来越少的能源,不具备自我恢复功能。
\end{solution}

% source: docs/review_questions.md | chapter=1. 绪论 | question=2
\begin{problem}\label{exer:think_01_intro_02}
什么是能源弹性系数，请给出其数学表达式？
\end{problem}

\begin{solution}
\begin{itemize}
\item \textbf{基本定义}：能源弹性系数是研究能源消费(生产)量与宏观经济发展指标(一般采用国内生产总值GDP)之间关系的数值。基本定义为能源消费(生产)增长率与经济增长率之比 。
\item \textbf{数学表达式}：
\end{itemize}
\begin{equation*}
e=\frac{\alpha}{\beta}=\frac{\Delta E/E_{0}}{\Delta G/G_{0}}
\end{equation*}
\begin{itemize}
\item \textbf{式中含义}：\(e\) 为能源消费(生产)弹性系数；\(\alpha\) 为能源消费(生产)增长率；\(\beta\) 为国内生产总值增长率；\(\Delta E\) 为计算期间的能源消费(生产)增量；\(E_{0}\) 为计算起始年的能源消费(生产)量；\(\Delta G\) 为计算期间的国内生产总值增量；\(G_{0}\) 为计算起始年的国内生产总值 。
\end{itemize}
\end{solution}

% source: docs/review_questions.md | chapter=1. 绪论 | question=3
\begin{problem}\label{exer:think_01_intro_03}
什么是二次能源，二次能源与一次能源相比有哪些优点？
\end{problem}

\begin{solution}
\begin{itemize}
\item \textbf{定义}：二次能源是指需要经过加工转换过程，从一次能源直接或间接转化而来的能源 。
\item \textbf{优点}：二次能源比一次能源具有更高的终端利用效率，使用时更方便、更清洁 。
\end{itemize}
\end{solution}

% source: docs/review_questions.md | chapter=1. 绪论 | question=4
\begin{problem}\label{exer:think_01_intro_04}
常见的能源环境污染分为哪几种？
\end{problem}

\begin{solution}
能源环境污染主要分为三大类，大体上可分为：\textbf{大气污染}、\textbf{水污染}以及\textbf{土壤污染} 。
\end{solution}

% source: docs/review_questions.md | chapter=1. 绪论 | question=5
\begin{problem}\label{exer:think_01_intro_05}
请简述改善能源环境污染的主要途径。
\end{problem}

\begin{solution}
\textbf{预防方面}：
\begin{enumerate}
\item 采取洗煤、脱硫脱硝等技术，降低污染物排放 ；
\item 实行汽车尾气排放标准（法律手段） ；
\item 汽车尾气处理 。
\end{enumerate}
\textbf{治理方面}：
\begin{enumerate}
\item 土壤修复技术、水污染治理 ；
\item 循环经济、环评制度 ；
\item 新能源技术：风能、太阳能、生物质能、海洋能、水能等各类新能源，在使用过程中几乎不产生对环境有害的物质 。
\end{enumerate}
\end{solution}

% source: docs/review_questions.md | chapter=1. 绪论 | question=6
\begin{problem}\label{exer:think_01_intro_06}
什么是温室效应，通常所说的温室气体主要有哪几种？
\end{problem}

\begin{solution}
\begin{itemize}
\item \textbf{温室效应}：由于大气中某些物质的增加，导致热量无法向外层空间发散，地球表面的温度则会逐渐升高。由于这种情况类似于温室大棚的保温效果，故称之为温室效应 。
\item \textbf{通常所说的温室气体种类}：1997年12月制定的《京都议定书》中有6种：二氧化碳（\(CO_2\)）、甲烷（\(CH_4\)）、氧化亚氮（\(N_2O\)）、氢氟碳化物（HFCs）、全氟化碳（PFCs）、六氟化硫（\(SF_6\)） 。
\end{itemize}
\end{solution}

% source: docs/review_questions.md | chapter=1. 绪论 | question=7
\begin{problem}\label{exer:think_01_intro_07}
碳排放因子是什么？有哪些二氧化碳减排措施？
\end{problem}

\begin{solution}
\begin{itemize}
\item \textbf{碳排放因子（排放因子 EF）}：单位某排放源使用量所释放的温室气体数量，单位与活动数据的单位相匹配 。
\item \textbf{二氧化碳减排措施}：
\begin{enumerate}
\item 提高能源利用效率 ；
\item 寻找清洁的新能源作为代替 ；
\item 二氧化碳的捕集与封存 。
\end{enumerate}
\end{itemize}
\end{solution}

% source: docs/review_questions.md | chapter=1. 绪论 | question=8
\begin{problem}\label{exer:think_01_intro_08}
举例说明你所熟悉的新能源发电技术？
\end{problem}

\begin{solution}
新能源发电技术是区别于传统化石能源的能源利用与转换的技术 。常见的技术类型有 ：
\begin{itemize}
\item \textbf{核能发电}：将原子核发生裂变或聚变时释放出的能量转化为电能 。目前商业核电站主要基于核裂变链式反应（如轻水反应堆、重水反应堆、石墨冷气堆） 。
\item \textbf{风力发电}：利用风力发电机组将风的动能转化为机械能再最终转化为电能，主要类型包括水平轴风力发电机和垂直轴风力发电机 。
\item \textbf{太阳能发电}：包括\textbf{太阳能光伏发电}（利用光生伏特效应原理将太阳光能直接转化为电能）和\textbf{太阳能光热发电}（将太阳能转化为中高温热能，再将热能转化为电能） 。
\item \textbf{生物质能发电}：相关技术包括直接燃烧发电、沼气发电和气化发电 。
\item \textbf{地热发电}：利用地下蒸汽或热水的热能，通过一定的装置和技术手段将其转化为机械能再转化为电能（包括双循环地热发电、闪蒸型地热发电、蒸汽型地热发电） 。
\item \textbf{海洋能发电}：包括潮汐能、波浪能、海流能、海洋温差能及盐差能等，江厦电站是中国最大的潮汐电站 。
\item \textbf{燃料电池发电}：一种把燃料所具有的化学能直接转换成电能的化学装置，如氢氧燃料电池 。
\end{itemize}
\end{solution}

% source: docs/review_questions.md | chapter=1. 绪论 | question=9
\begin{problem}\label{exer:think_01_intro_09}
新能源渗透率如何定义？
\end{problem}

\begin{solution}
\textbf{新能源渗透率}：各种新能源占全部能源用量的比重 。
\end{solution}

% source: docs/review_questions.md | chapter=1. 绪论 | question=10
\begin{problem}\label{exer:think_01_intro_10}
可以试着描述一下风光互补发电系统的适用领域及其优势吗？
\end{problem}

\begin{solution}
\textbf{适用领域}：
\begin{enumerate}
\item \textbf{无电及偏远地区}：如远离电网的孤岛、边防哨所、偏远山区及牧区，用于解决基本生活供电。
\item \textbf{交通与公共基础设施}：如风光互补路灯、城市景观照明、高速公路监控及野外气象观测站。
\item \textbf{通信基站}：为建立在山顶或偏远地区的移动通信基站提供高可靠性的独立电源。
\item \textbf{农业与生态应用}：野外农业提水灌溉、温室大棚监测设备及生态园区的微电网。
\end{enumerate}
\textbf{核心优势}
\begin{enumerate}
\item \textbf{资源互补性强}：利用“白天光照好、夜间风大”以及“夏季光强风弱、冬春风大光弱”的自然规律，实现昼夜与季节互补，发电更连续。
\item \textbf{供电稳定性高}：双能源输入减少了单一能源受天气限制（如连续阴雨或无风）带来的断电风险，系统可靠性更高。
\item \textbf{降低储能成本}：由于发电出力更均匀，系统对蓄电池容量的需求显著降低，从而减少了昂贵的电池投资和维护成本。
\item \textbf{节约资源与土地}：风机和光伏组件可空间复合布局，节约占地；同时两者共用一套逆变与输配电系统，有效降低整体建设成本。
\end{enumerate}
\end{solution}

\section{计算题}

% index: calc_01_intro_p041_01 | source=others/01_绪论.pdf | page=41 | slide=1.1.3 能源结构-能源弹性系数 | status=checked
\begin{problem}[能源消费或生产弹性系数]\label{exer:calc_01_intro_p041_01}
已知计算期间能源消费或生产增量为 \(\Delta E\)，计算起始年的能源消费或生产量为 \(E_0\)；同期国内生产总值增量为 \(\Delta G\)，计算起始年的国内生产总值为 \(G_0\)。写出能源消费或生产弹性系数 \(e\) 的计算式，并说明各量含义。
\end{problem}

\begin{solution}
基本定义为能源消费或生产增长率与经济增长率之比：
\begin{equation}
e = \frac{\alpha}{\beta}
  = \frac{\Delta E / E_0}{\Delta G / G_0}.
\end{equation}
式中，\(e\) 为能源消费(生产)弹性系数；\(\alpha\) 为能源消费(生产)增长率；\(\beta\) 为国内生产总值增长率；\(\Delta E\) 为计算期间的能源消费(生产)增量；\(E_{0}\) 为计算起始年的能源消费(生产)量；\(\Delta G\) 为计算期间的国内生产总值增量；\(G_{0}\) 为计算起始年的国内生产总值 。
\end{solution}

% index: calc_01_intro_p087_01 | source=others/01_绪论.pdf | page=87 | slide=1.3.1 温室效应 | status=checked
\begin{problem}[温室气体全球变暖潜能值查表]\label{exer:calc_01_intro_p087_01}
根据温室气体全球变暖潜能值表，列出 \(100\) 年时间尺度下 \(CO_2\)、\(CH_4\)、\(N_2O\)、\(CF_4\)、\(C_2F_6\)、\(SF_6\) 的 GWP 值，并说明如何由温室气体质量换算为二氧化碳当量。
\end{problem}

\begin{solution}
\(100\) 年全球变暖潜能值分别为：
\begin{equation*}
\begin{aligned}
GWP_{CO_2} &= 1, &
GWP_{CH_4} &= 21, &
GWP_{N_2O} &= 310,\\
GWP_{CF_4} &= 6500, &
GWP_{C_2F_6} &= 9200, &
GWP_{SF_6} &= 23900.
\end{aligned}
\end{equation*}
若某温室气体质量为 \(m\)，其二氧化碳当量按
\begin{equation}
m_{CO_2e} = m \cdot GWP
\end{equation}
换算。
\end{solution}

% index: calc_01_intro_p108_01 | source=others/01_绪论.pdf | page=108 | slide=1.3.4 碳核算 | status=checked
\begin{problem}[质量平衡法碳核算]\label{exer:calc_01_intro_p108_01}
某排放源采用质量平衡法核算温室气体排放量。设输入物料量为 \(M_I\)，输出物料量为 \(M_O\)，输入和输出物料含碳量分别为 \(CC_I\)、\(CC_O\)，碳质量转化为温室气体质量的转换系数为 \(w\)，全球变暖潜能值为 \(GWP\)。写出温室气体排放量 \(E_{GHG}\) 的计算式。
\end{problem}

\begin{solution}
根据质量守恒，用输入物料含碳量减去输出物料含碳量，再乘以转换系数和全球变暖潜能值：
\begin{equation}
E_{GHG}
= \left[\sum(M_I \cdot CC_I) - \sum(M_O \cdot CC_O)\right]\cdot w \cdot GWP.
\end{equation}
式中，\(E_{GHG}\) 的单位可写作 \(\unit{tCO_2e}\)。该页未给具体物料数据，不能进一步数值化。
\end{solution}
